Differential testing rests on a reference the system under test cannot influence.
In practice, ``independence'' is almost always discussed along a single axis: how
independently the reference was \emph{implemented}---by different authors, from a
different specification, in a separate codebase. A recent systematic review of
oracle sources grades precisely this, ranking an independently developed reference
above cross-implementation comparison, and that above consensus among
implementations drawn from a common origin~\cite{oraclesurvey2026}. A second axis
is rarely named: what the reference \emph{consumes}. An oracle whose inputs are
artifacts produced by the system under test inherits that system's defects, no
matter how independently its own logic was written.

The clearest sign that the second axis is unnamed rather than unimportant is a
technique that already depends on it. Csmith generates random C programs and compares
what several compilers make of them, and its authors meet the independence question
head on: they state that a defect shared by every compiler would go undetected, and
answer the concern empirically, reporting no correlated failures among unrelated
compilers~\cite{yang2011csmith}. That answer is about how independently the compilers
were built. A second reason is available and the paper does not give it: the program
under test is produced by Csmith, hence by none of the compilers being compared, so
the evidence never passes through a system under test. Csmith is protected by an axis
it does not name, and defends itself on the one it does. The axis was load-bearing
before it was articulated.

This is not a fine distinction that practice already respects; practice pushes the
other way. Artifacts produced by the system under test---intermediate
representations, execution profiles, planner estimates, internal state dumps---are
the richest and cheapest signal available, and oracles are built on them
\emph{because} they are rich. Equivalence Modulo Inputs derives its variants from a
profile of the program under test~\cite{le2014emi}; NoREC evaluates its reference on
the very engine it is testing~\cite{rigger2020norec}. Both are effective, and both
are evidentially entangled with some component of the system they judge. The
convenient design is the one that violates the axis.

The consequence is not vague. If an oracle inherits defects in what it consumes,
its blind-spot class is determined before the oracle exists: partition the system's
historical defect population by whether each defect corrupts an artifact the oracle
would consume, and the partition bounds the oracle's reach. This costs source
reading---no builds, no execution, no test campaign. Yet it cannot be assembled from
published results. Defect taxonomies in the testing literature are organized by
symptom (wrong code, crash, performance) and by configuration (32- against 64-bit
targets), never by whether an oracle could have observed the defect. Where a technique
does state what it cannot see---and some state it precisely---it states it \emph{by
component}: which part of the system was analysed. That is a different partition, and
not convertible into this one, because two oracles covering the same component can
differ entirely in what they consume.

We name the axis, give the procedure that operationalizes it, and apply it once and
deeply. The subject is the Linux eBPF verifier, chosen because it is already heavily
fuzzed, because its defects are individually documented together with the state they
corrupt, and because its correctness is safety-relevant. Two oracles were built for
it under identical discipline---same corpus, same author, same
refusal-over-approximation rule, both independent reimplementations---differing in
one respect that is visible in their type signatures: one consumes the emitted
instruction stream, the other consumes the state the verifier itself wrote. Their
blind-spot classes separate exactly as the axis predicts, and each side is anchored
by calibration against real defects: sixteen pairs, most built from a kernel that
carried the defect and one that did not, the rest settled from the defective state the
fix's own message prints. Applying the procedure ahead of time to
fifty-seven historical fixes---those whose messages quote internal state, a frame we
later measure to be enriched about threefold for the oracle's own path---shows that
the more sophisticated of the two oracles is structurally blind to 42\% of that
population and can target 2\% of it; on the unfiltered complement both shares fall.
Either way the result was available before the oracle was written.

A differential oracle's blind-spot class follows from where its reference reads, not
from how independently it was written; and because it follows, it can be measured
before the oracle is built rather than discovered after it has failed to fire.
